\chapter{Visual Paradigm Modelling and Reproduction Guide}
\section{Native project and workflow}
The editable project is \texttt{diagrams/AtelierFashion.vpp}, created in Visual Paradigm 18 Community Edition. Every report diagram is composed of native UML shapes and relationships in that application. The local Open API plugin creates the models, arranges their bounds, saves the project, and invokes Visual Paradigm's own PNG export at scale 2. No diagram is substituted by PlantUML, Mermaid, drawing-library output, or an externally generated raster. Colour, font, and connector-style settings are left at the application's defaults; only positions, dimensions, caption placement, and necessary UML semantics are adjusted. The Community Edition watermark is retained.

For manual editing, open the VPP file, select the named diagram in Diagram Navigator, edit model properties and shape bounds, and save. Export with Project $>$ Export $>$ Diagrams as Image, choosing PNG. Inspect the exported image as well as the in-application canvas: labels must be complete, nodes must not overlap, connectors must reach their intended endpoints, and unrelated connectors/labels must not collide. Review the inserted image again at PDF page size. Wide class and interaction diagrams use landscape report pages for legibility.

\section{Purpose, content, and construction by diagram}
\begin{longtable}{P{.15\linewidth}P{.3\linewidth}P{.43\linewidth}}
\caption{Diagram construction and review guide}\label{tab:vp-guide}\\
\toprule Export & Purpose and content & Visual Paradigm construction/review procedure \\\midrule\endfirsthead
\toprule Export & Purpose and content & Visual Paradigm construction/review procedure \\\midrule\endhead
01 customer use cases & Customer goals and payment boundary & New Use Case Diagram; add System, Visitor, Customer, Payment Provider and ellipses; place ellipses inside the system. Add Customer generalisation toward Visitor, actor associations, and Place order include arrows toward validation, reservation, and payment initiation. Inspect all arrowheads and captions.\\
02 operations & Catalogue, stock, fulfilment, inspection, carrier events & Add staff/provider actors outside the same system boundary and UC09--UC13 ellipses inside. Associate the carrier with shipment support and its event goal. Route the provider link around adjacent use cases and check that every line touches its ellipse.\\
16 back office & Cancellation, returns, refunds, reports, access & Use Case Diagram with Customer Service, Finance, Manager, Administrator and Payment Provider. Associate each role with its permitted goals. Check that the finance/report fan does not cross unrelated ellipses.\\
18 supporting goals & Notifications, optional reviews, financial reconciliation & Reuse Customer, Payment Provider, and UC05; add UC17 and UC18 and Finance Staff's reconciliation participation. Keep actors outside the system and distinguish the optional review feature in the requirements.\\
03 search & Variant filters and unavailable/no-match retries & New Activity Diagram; use Initial Node, Action, Decision Node, Merge Node and Activity Final Node. Place mutually exclusive guards on decision outputs. Route retry paths to explicit merges before re-entering actions. Verify both labels and control semantics.\\
04 checkout & Customer, system, and provider responsibilities & Activity Diagram with three vertical swimlane partitions. Put consent in Customer, local allocation in System, hosted payment in Provider, and verified reconciliation back in System. Review lane ownership and committed/rejected branches.\\
05 inspection & Physical receipt before stock restoration & Three partitions: Customer Service, System, Warehouse. Place approval-reference check, physical inspection, receipt validation, and once-only restock in their responsible lanes. Give invalid/valid receipt branches separate visible destinations.\\
14 receiving & Validated, auditable stock receipt & Activity Diagram with receipt identification, validation, guarded posting, duplicate/invalid alternative and final nodes. Confirm the invalid branch never passes through the stock-posting action.\\
15 remedy & Refund versus reserved exchange & Activity Diagram with remedy and stock decisions. Add explicit refund, exchange, allocated/unavailable guards. Link successful allocation to replacement fulfilment; unavailable stock ends the attempt awaiting customer choice.\\
06 order states & Permitted fulfilment lifecycle & New State Machine Diagram; use Initial Pseudostate, State and Transition. Add PendingPayment, Confirmed, Packing, Shipped, Delivered, Cancelled, Expired and DeliveryFailed. Use short event captions and the report transition table for guards/effects. Route expiry and permanent failure outside state boxes.\\
07 catalogue classes & Category, Product, Variant, SizeChart, InventoryBalance & New Class Diagram; add representative attributes and associations with multiplicities. Keep product/variant distinct and optional chart use explicit. Space multiplicity labels away from class corners.\\
08 shopping/order classes & Customer, Cart, CartLine, Order, OrderLine, AddressSnapshot and Variant & Reuse existing model classes when adding views. Associate cart lines with exact SKU and committed orders with snapshots. Route Order--OrderLine outside the Variant box; place AddressSnapshot where its link avoids the cart/SKU connector.\\
17 reservations & Explicit commitment separate from carts & Class Diagram reusing Order, Variant and InventoryBalance; add Reservation with quantity, kind, state and deadline. Associate each reservation with one original order and one SKU. Check that diagonal associations do not intersect the inventory link.\\
09 after-sales classes & Payment, Refund, Shipment, ReturnRequest, ReturnLine and ExchangeFulfilment & Reuse Order and OrderLine. Show one payment per refund, optional original shipment, per-line returns and optional exchange fulfilment. Interpret quantity/money caps through rules, not multiplicity alone.\\
10 checkout sequence & Commit local work before provider interaction & New Sequence Diagram; add customer, checkout service, transactional store and provider lifelines. Use synchronous calls and dashed replies in vertical time order. Steps 2--5 share a local transaction; no database lock extends into the provider call.\\
11 callback sequence & Durable receipt, guarded update, optional compensation & Add provider, webhook adapter, payment service and store lifelines. Acknowledge persisted receipt; use an asynchronous worker message; add an opt fragment guarded by closed and paid for compensation. Check guard text against lifelines and message labels.\\
12 components & Application and integration dependencies & New Component Diagram; add UI, application, repository, worker, adapters and external services; add directed dependencies. Use concise labels that fit gaps. Verify actual component names after saving, then export.\\
13 deployment & HTTPS entry, private data access, backup, provider callbacks & New Deployment Diagram; add native Node shapes and communication associations. Route callbacks around other nodes and connections; keep all link labels within export bounds.\\
\bottomrule
\end{longtable}

\section{Local automation and report rebuild}
The modelling source is in \texttt{automation/FashionPlugin.java}; diagram specifications and generators are in \texttt{automation/}. The plugin runs inside Visual Paradigm and consumes a tab-separated job file. Completion and errors are written separately; a generator exit is not itself proof that Visual Paradigm completed a diagram. After each completed export, the image is reviewed and any defects are corrected before the next model is accepted. The native model audit records diagram/model identifiers and associations independently of the PNGs.

The report source is \texttt{report/main.tex} with separate chapter files. From the report directory, run the locally provided Tectonic executable:
\begin{verbatim}
..\automation\tectonic\tectonic.exe --keep-logs main.tex
\end{verbatim}
Tectonic retrieves required LaTeX packages on its first run and caches them. The source archive excludes tool binaries; install Tectonic and run \texttt{tectonic main.tex} when rebuilding elsewhere. A standard TeX installation may alternatively build the same source with repeated pdfLaTeX runs until references stabilise. Keep the sibling \texttt{diagrams} folder because figures are included by relative paths. The final review checks compilation messages, reference resolution, page layout, and actual rendered figures; the report does not equate this artifact check with implementation testing.

\chapter{Source Use and Assumption Register}
\section{Use of the supplied examples}
The supplied files are retained unchanged in \texttt{third\_party}. Their section structures informed coverage, while their domain choices were assessed against the retailer boundary. The fashion example suggested variant, sizing, and after-sales concerns; the marketplace example highlighted the need to distinguish stock commitments, transaction status, and fixed order evidence. Independent seller payouts and perishable expiry management were excluded. Wording, student identity, and diagrams from those reports were not adopted as the final report's own artifacts. They are local examples, not authoritative standards.

\section{Theory and commercial references}
OMG UML 2.5.1 supplies modelling notation. Shopify's public product/variant model supports the design/SKU distinction as an industry example. Stripe's payment, webhook, and idempotency documentation informs the handling of uncertain outcomes and repeated/out-of-order messages. OWASP ASVS and W3C WCAG 2.2 define review references for security and accessibility targets. Visual Paradigm's own Open API documentation supports native model creation and export. These references are cited where their ideas are used; the proposed policies and architecture remain the author's case-design decisions.

\section{Assumptions to confirm}
One retailer/warehouse/currency, authenticated checkout, online prepaid payment, one original shipment, tax-inclusive listed prices, 15-minute initial holds, the 14-day return request window, disclosed return exclusions, equal-price same-product exchanges, shipping refund policy, and numerical service targets are case assumptions. The study does not establish legal eligibility, provider market availability, realised business improvement, production performance, or completed stakeholder interviews. Validation with the actual retailer must confirm or revise these assumptions before implementation.
